% ============================================================
% ANEXO B - Respuestas preparadas (Kamila Calle - PE5)
% Bloques asignados: Inteligencia artificial (16-18) y Etica (21-22)
% Fragmento para \input en el informe final
% ============================================================
\subsection*{Inteligencia artificial}

\noindent\textbf{16. ¿Qué decide el modelo y qué pasa cuando se equivoca?}\\
Ningún componente decide: ambos sugieren. IA-01 genera recomendaciones (RF-15) y alertas tempranas (RF-26) con un nivel de confianza explícito (D-03); IA-02 propone una clase del catálogo de plagas y enfermedades (RF-20, D-04) a partir de una fotografía. Ninguno actúa solo: la restricción R-05 exige confirmación humana, y esto es verificable en dos requisitos concretos: RF-28 (ninguna alerta o recomendación se aplica sin que el administrador la confirme o descarte) y RF-30 (una propuesta de clasificación descartada no crea ni modifica ningún registro). Si el modelo se equivoca, el costo es una verificación en campo innecesaria o una alerta omitida, nunca una acción automática errónea. \emph{Fallback}: si el módulo de IA-01 falla, el resto del sistema sigue operando sin recomendaciones ni alertas (RNF-18); si IA-02 no está disponible o no hay conexión, el registro continúa por selección manual del catálogo (RNF-23). \textbf{Artefactos:} fichas de componente (ERS v2.0, Sección 3.4), tabla de decisiones D-01 a D-05, RNF-18 y RNF-23.

\medskip
\noindent\textbf{17. ¿Con qué datos lo entrenarían y con qué base legal?}\\
IA-01 se entrena con registros propios del sistema (producción, actividades, enfermedades, inventario), con un mínimo operativo de 500 registros por ciclo anual; por debajo de ese volumen opera solo con reglas, sin capa predictiva. IA-02 se entrena con entre 800 y 1200 fotografías propias de las parcelas (mínimo 150 por clase), complementadas con conjuntos públicos como semilla inicial, con doble etiquetado y concordancia mínima kappa de 0,80 entre etiquetadores. Ambos bloques de datos son aporte de María Escudero, documentados en la ficha de cada componente. Base legal LOPDP: los registros que entrenan IA-01 contienen datos personales del trabajador agrícola (identificador, actividad, parcela); el tratamiento se ampara en la ejecución de la relación laboral y el interés legítimo de gestión agrícola (Art.~8 LOPDP), con finalidad declarada de planificación productiva, no de evaluación individual. Las fotografías que entrenan IA-02 no son datos personales bajo el protocolo de captura (exclusión de personas del encuadre), por lo que este componente no añade tratamiento de datos personales adicional. \textbf{Artefactos:} bloque de datos de entrenamiento de cada ficha de componente (ERS v2.0, Sección 3.4).

\medskip
\noindent\textbf{18. ¿Cómo medirían que el modelo no perjudica sistemáticamente a un grupo?}\\
Con los cuatro RNF de equidad del catálogo, uno por cada eje de riesgo identificado en el propio sistema: RNF-19 limita a 10 puntos porcentuales la brecha de exactitud de IA-01 entre los cultivos de cacao y verde (tolerancia mayor porque son dominios agronómicos distintos); RNF-20 limita a 5 puntos porcentuales la brecha de falsos positivos entre parcelas con alto y bajo volumen histórico de registro, porque las parcelas revisadas con más frecuencia acumulan más incidencias y podrían sesgar el modelo a su favor; RNF-24 y RNF-25 limitan a 5 puntos porcentuales, respectivamente, la brecha de F1 de IA-02 entre condición de luz de la fotografía y entre gama de dispositivo de captura, los dos sesgos de datos que el propio equipo documentó al preparar el conjunto de imágenes. El método es el mismo en los cuatro casos: calcular la métrica de desempeño por separado para cada grupo y comparar contra el umbral; si lo supera, se trata como un defecto del modelo, no como una particularidad del grupo, y se amplía la recolección de datos del grupo perjudicado antes de reentrenar. \textbf{Artefactos:} RNF-19, RNF-20, RNF-24 y RNF-25 con su método de medición; bloque de equidad de cada ficha de componente.

\subsection*{Ética}

\noindent\textbf{21. ¿Qué datos personales trata el sistema y cómo se protegen?}\\
El SGA trata identidad y actividades de los trabajadores agrícolas (RF-14) y la autoría de los registros y reportes que generan (RF-04, RF-05, RF-16). La protección combina tres mecanismos ya especificados en el catálogo: control de acceso diferenciado por rol (RF-01), cifrado de las credenciales almacenadas (RNF-11) y, para los datos que alimentan IA-01, una base legal explícita (ejecución de la relación laboral e interés legítimo de gestión agrícola, Art.~8 LOPDP) con una finalidad declarada que excluye la evaluación individual del trabajador. IA-02 no añade riesgo nuevo sobre datos personales: opera sobre fotografías de plantas, y el protocolo de captura excluye personas, documentos o placas del encuadre. \textbf{Artefactos:} RF-01, RF-14, RNF-11; bloque de datos de entrenamiento y clasificación de riesgo de cada ficha de componente (ERS v2.0, Sección 3.4).

\medskip
\noindent\textbf{22. ¿Qué partes del informe fueron asistidas por IA y cómo las validaron?}\\
Se responde con el Anexo E de este informe, que declara por sección la herramienta usada, el tipo de asistencia y el método de validación: en las secciones a mi cargo (Sec.~3 ERS/SRS final, Sec.~7 Requisitos de IA, y las fichas RF-28 a RF-31 y RNF-16 a RNF-26 del ERS v2.0) la asistencia fue de redacción y estructuración a partir de la evidencia ya citada en la línea base del ERS y de decisiones tomadas por mí, nunca de generación de criterio propio sobre qué elegir o qué umbral fijar. Cada fuente citada (EV-01, EV-02) se verificó contra el propio texto del ERS, que ya trae esas citas para los requisitos existentes a los que se enlazan los nuevos; la clasificación de riesgo se construyó aplicando los criterios del Reglamento (UE) 2024/1689 al sistema real, no copiando un ejemplo genérico. Las decisiones evaluativas (qué componente de IA elegir, qué clasificación de riesgo aplica, qué umbral fija cada RNF) son producción propia. \textbf{Artefacto:} Anexo E (declaración de uso de IA), firmado por los cinco integrantes.
